\subsection{Arithmetic detection and the limits of fixed operation invariants}

\begin{proposition}[C33: \status{Proved}; bounded germs and fixed operation groups]
\label{sa08:C33}
Let $S=\Gamma_{\Z}(T)$, let $\mathcal B$ be the Serre subcategory
of nonnegatively graded $S$-modules with bounded weight support,
and let $\mathcal T$ be its localizing Serre closure.
\begin{enumerate}
\item Arithmetic localization does not detect zero objects of the
bounded-weight quotient $\operatorname{Gr}_{\geq0}(S)/\mathcal B$.
There is a nonzero object whose rational and every $p$-local
bounded-weight germs vanish.
\item If a graded $S$-module $M$ has bounded rational support and
bounded $p$-local support for every prime $p$, then $M\in\mathcal T$.
\item Let $G$ be one fixed finitely generated abelian group and
$g_w\in G$. If, for every prime $p$, the image of $g_w$ in
$G\otimes\Z_{(p)}$ is zero for all sufficiently large $w$, then
$g_w=0$ for all sufficiently large $w$.
\end{enumerate}
The bounds in the hypotheses may depend on the prime.
\end{proposition}
\begin{proof}
For the first assertion choose pairwise distinct primes $p_w$,
one for each positive weight, and set
\[
 M_w=\Z/p_w\quad(w\geq1),\qquad M_0=0,\qquad S^{>0}M=0.
\]
This is a graded $S$-module with unbounded support, so its image
in the Serre quotient is nonzero. Rationalization is zero.
For any fixed $p$, localization at $p$ leaves at most one nonzero
weight, so that localized germ is zero. This module is not
asserted to be a generalized theory or an invariant of two theories.

For the second assertion choose $N$ with
$M_{\geq N}\otimes\Q=0$. The tail is an $S$-submodule and a
torsion abelian group. Its canonical primary decomposition
\[
 M_{\geq N}=\bigoplus_p (M_{\geq N})[p^\infty]
\]
is preserved by $S$. Each summand is its own $p$-localization,
and therefore has bounded weight support by hypothesis.
Closure under arbitrary direct sums puts the tail in
$\mathcal T$. The quotient by the tail is bounded, and closure
under extensions gives $M\in\mathcal T$. In particular the
example in the first assertion is erased by this larger quotient.

For the third assertion let $\mathcal P$ be the finite set of
torsion primes of $G$ and add one arbitrary prime $p_0$.
The homomorphism
\[
 G\longrightarrow
 \bigoplus_{p\in\mathcal P\cup\{p_0\}}G\otimes\Z_{(p)}
\]
is injective: the added prime detects the free part and each
primary torsion part is detected at its own prime. Taking the
maximum of the finitely many eventual bounds proves the assertion.
\end{proof}

\paragraph{Conditional application to ordinary primary operations.}
Fix positive input degrees $i_1,\ldots,i_r$ and an output degree $j$.
Every natural operation on simplicial sets
\[
 \prod_{a=1}^r H^{i_a}(-;\Z)\longrightarrow H^j(-;\Z)
\]
is determined by its value on the universal tuple on
$K=\prod_a K(\Z,i_a)$. Indeed every input tuple is represented
by a map to $K$, and naturality recovers its value by pullback.
No additivity assumption is needed. The group
$G=H^j(K;\Z)$ is finitely generated: Serre's finite-generation
class and his theorem on Eilenberg--MacLane homology give finite
generation in each homological degree; finite-product K\"unneth
and the universal coefficient sequences give the assertion
for cohomology\cite[Chapter I, \S6.1, p.~264; Chapter II,
\S6, Proposition~8, p.~271]{BSerre1953Classes}.
The same sequences identify
$G\otimes\Z_{(p)}$ with $H^j(K;\Z_{(p)})$.

Suppose two families of weighted operations give operations of
this fixed ordinary type, commute with coefficient localization,
and are respected by comparison maps of theories. If the local
comparison theorems identify their classes eventually at each
prime, the differences of their universal values satisfy
C33(3). Thus their integral operation germs agree after one
weight bound. This conditional argument does not assume that
every cochain divided power descends to cohomology.

\paragraph{Scope for the unresolved integral problem.}
The first part exhibits the nonuniform-prime phenomenon that an
actual eventual construction must control. The second part shows
why the larger localizing shadow quotient can forget that
information. The third part prevents one fixed finitely generated
ordinary operation group from detecting such a difference.
Allowing the operation degree or the coefficient prime to vary
does not satisfy that fixed-group hypothesis. None of these
assertions constructs coherent comparison maps, an admissible
counterexample, or the local repairs required by
Statements~\ref{sa08:C30} and~\ref{sa08:C32}.
